%% ITEM 0 TYPE Section
\section{Benchmark results and efficiency}%% ITEM 1 TYPE Section
\section{Benchmark Results and Efficiency}%% ITEM 2 TYPE Frame
\begin{frame}
\deepslideid{S0001}

\frametitle{Benchmark Overview}

\begin{itemize}
\item Comprehensive evaluation across language and vision-language tasks
\item Strong baselines: GPT5.2, Claude 4.5 Opus, Gemini-3 Pro
\item Key categories:
  \begin{itemize}
  \item Language: Knowledge, reasoning, agentic tasks, coding
  \item Vision-Language: STEM/puzzle, VQA, spatial intelligence, video
  \end{itemize}
\item Qwen3.5-397B-A17B shows competitive performance
\end{itemize}
\end{frame}%% ITEM 3 TYPE Frame
\begin{frame}
\deepslideid{S0002}

\frametitle{Language Benchmark Performance}

\begin{itemize}
\item MMLU: 89.5 (vs GPT5.2: 90.1, Claude 4.5: 89.8)
\item Big-Bench Hard: 89.2 (competitive with SOTA)
\item GPQA (STEM reasoning): 62.1
\item HumanEval (coding): 90.2
\item MBPP (coding): 86.7
\item Strong performance across reasoning, knowledge, coding
\end{itemize}
\end{frame}%% ITEM 4 TYPE Frame
\begin{frame}
\deepslideid{S0003}

\frametitle{Vision-Language Benchmark Results}

\begin{itemize}
\item STEM/puzzle: Strong performance on science/math reasoning
\item General VQA: Competitive on visual question answering
\item Text/document: Excellent OCR and document understanding
\item Spatial intelligence: Good geometric/spatial reasoning
\item Video understanding: Competitive temporal reasoning
\item Medical VQA: Specialized medical image analysis
\end{itemize}
\end{frame}%% ITEM 5 TYPE Frame
\begin{frame}
\deepslideid{S0004}

\frametitle{Efficiency Metrics: Throughput and Latency}

\begin{figure}[htbp]
  \centering
  \includegraphics[width=0.85\linewidth]{../qwen35\_paper/figures/decode\_throughput.png}
  \caption{Decode throughput comparison across models and configurations.\cite{gigazine_qwen35}}
\end{figure}

\begin{itemize}
\item High throughput with optimized inference
\item Low latency for real-time applications
\item Efficient parameter-activation scaling via MoE
\end{itemize}
\end{frame}%% ITEM 6 TYPE Frame
\begin{frame}
\deepslideid{S0005}

\frametitle{Parameter-Activation Scaling via MoE}

\begin{itemize}
\item Mixture of Experts architecture enables efficient scaling
\item Only relevant experts activated per token
\item Reduces computational cost while maintaining capacity
\item Enables larger model sizes with manageable inference cost
\item Supports diverse task capabilities through specialized experts
\item Efficient memory usage during inference
\end{itemize}
\end{frame}%% ITEM 7 TYPE Frame
\begin{frame}
\deepslideid{S0006}

\frametitle{Summary and Key Takeaways}

\begin{itemize}
\item Competitive performance across language and vision-language benchmarks
\item Strong agentic/tool-use capabilities
\item Efficient inference with high throughput and low latency
\item MoE architecture enables scalable parameter-activation efficiency
\item Practical for real-world deployment scenarios
\item Open-source availability enhances accessibility
\end{itemize}
\end{frame}%% ITEM 8 TYPE Frame
\begin{frame}
\deepslideid{S0007}

\frametitle{Conclusion}

\begin{itemize}
\item Qwen3.5-397B-A17B delivers strong benchmark performance
\item Competitive with leading proprietary models
\item Efficient architecture enables practical deployment
\item Open-source availability benefits research community
\item Future work: further optimization and specialized applications
\end{itemize}
\end{frame}%% ITEM 9 TYPE Section
\section{Improvements and limitations}%% ITEM 10 TYPE Section
\section{Improvements and Limitations}%% ITEM 11 TYPE Frame
\begin{frame}
\deepslideid{S0008}

\frametitle{Key Improvements in Qwen3.5}

\begin{itemize}
    \item \textbf{Enhanced multimodal agentic performance}
    \begin{itemize}
        \item Early-fusion training on large multimodal corpora
        \item Better alignment of visual understanding with language reasoning
        \item Improved performance on vision-language benchmarks
    \end{itemize}
    
    \item \textbf{Computational efficiency gains}
    \begin{itemize}
        \item Gated Delta Networks for reduced inference overhead
        \item Sparse MoE routing (512 experts, 10+1 activated)
        \item 397B total parameters with only 17B active
    \end{itemize}
    
    \item \textbf{Scalable RL generalization}
    \begin{itemize}
        \item RL scaled across large collections of tool-using environments
        \item Progressive complexity in task distributions
    \end{itemize}
\end{itemize}
\end{frame}%% ITEM 12 TYPE Frame
\begin{frame}
\deepslideid{S0009}

\frametitle{RL Environment Scaling}

\begin{figure}[htbp]
  \centering
  \includegraphics[width=0.85\linewidth]{../qwen35\_paper/figures/env\_scaling.png}
  \caption{RL environment scaling across progressively complex task distributions}
  \label{fig:env_scaling}
\end{figure}

\begin{itemize}
    \item Progressive scaling of RL training environments
    \item Increasing complexity in task distributions
    \item Supports tool-using agent capabilities
\end{itemize}
\end{frame}%% ITEM 13 TYPE Frame
\begin{frame}
\deepslideid{S0010}

\frametitle{Key Limitations and Open Questions}

\begin{itemize}
    \item \textbf{Unclear generalization bounds}
    \begin{itemize}
        \item Limited theoretical analysis of generalization capabilities
        \item Performance boundaries across diverse real-world scenarios not well-defined
        \item Transfer learning limits to novel domains uncertain
    \end{itemize}
    
    \item \textbf{Scaling behavior details insufficient}
    \begin{itemize}
        \item Exact scaling laws for multimodal performance not provided
        \item Efficiency trade-offs at different scales not quantified
        \item Optimal parameter allocation strategies unclear
    \end{itemize}
    
    \item \textbf{Architectural transparency gaps}
    \begin{itemize}
        \item Full architectural details of hybrid components not disclosed
        \item Training curriculum specifics for RL scaling limited
        \item Multimodal fusion mechanisms partially described
    \end{itemize}
\end{itemize}
\end{frame}%% ITEM 14 TYPE Frame
\begin{frame}
\deepslideid{S0011}

\frametitle{Architecture Efficiency Features}

\begin{itemize}
    \item \textbf{Sparse MoE design:} 512 experts, only 10+1 activated per token
    \item \textbf{Parameter efficiency:} 397B total vs 17B active parameters
    \item \textbf{Hybrid backbone:} Combines Gated DeltaNet and Attention layers
    \item \textbf{Extended context:} 262K native, extensible to 1M+ tokens
\end{itemize}

\vspace{0.5cm}
\textbf{Key trade-offs:}
\begin{itemize}
    \item Efficiency gains from sparse activation
    \item Complexity in expert routing and load balancing
    \item Memory overhead for expert parameters
\end{itemize}
\end{frame}